\appendix

\section{Sampling Bounds}\label{sec:sampling}

Recall the Chernoff bound:


\begin{proposition}\label{lem:chernoff}
    Let $X_1, \cdots, X_n$ be independent random variables with values in $[0,1]$ and $X = \sum_i X_i$. Then: \begin{itemize}
        \item \textbf{Additive form: } For any $b \ge 0$, $\Pr[|X - \E[X]| \ge b] \le 2\cdot \exp\left(-\frac{b^2}{2n}\right)$, \label{itm:chernoff-add}

        \item \textbf{Multiplicative form: } For $0 \le \delta \le 1$, $\Pr[|X - \E[X]| \ge (1+\delta) E[X]] \le \exp\left(-\frac{\delta^2 E[X]}{3}\right)$.\label{itm:chernoff-mult}
    \end{itemize}
\end{proposition}

\bigskip\noindent
Now suppose we are allowed to repeatedly sample from an unknown Bernoulli distribution $\mu$. We know that $\mu$ either has acceptance probability $p$ or acceptance probability $q$. We wish to determine which is the case, via some method that may err with probability $\gamma$. Suppose that, for some $1 \ge \alpha,\beta > 0$, $p \le \alpha$, whereas $q \ge (1 + \beta) \alpha$. We may then draw 
\[
k = k(\alpha, \beta, \gamma) = \left\lceil\frac{48}{\alpha \beta^2} \ln\left(\frac{1}{\gamma}\right)\right\rceil
\]
samples $X_1, \ldots, X_k$ from $\mu$, and choose the threshold
\[
\tau = \tau(\alpha, \beta, \gamma) = \alpha \cdot \left(1 + \frac{\beta}{4}\right) \cdot k 
\]
It now follows, from the multiplicative form of the Chernoff bound (\Cref{lem:chernoff}), that
\[
\Pr[\sum_i X_i \ge \tau] = \Pr[\sum_i X_i \ge \alpha k \;\left(1 + \frac\beta4\right)] \le \exp\left(- \frac{\beta^2}{48} \alpha k \right) \le \gamma.
\]
Now, whenever $\beta \in [0,1]$, we have $1 + \frac\beta4 \le (1 + \beta) ( 1 - \frac{\beta}{4})$, and so, if the acceptance probability is $q \ge \alpha(1 + \beta)$, then
\begin{align*}
\Pr[\sum_i X_i \le \tau] %
& = \Pr[\sum_i X_i \le \alpha k\left(1 + \frac\beta4\right)] \le \Pr[\sum_i X_i \le \alpha (1 + \beta) k\left(1 - \frac\beta4\right)] \\
& \le \exp\left(- \frac{\beta^2}{48} \alpha (1+\beta) k \right) \le \gamma.
\end{align*}





\section{Proof of Propositions Regarding Triangular Discrimination} \label{sec:td-props}

\lambdatvd*

\begin{proof}
To see the lower bound on $\Lambda(\mu, \nu)$:
\[
\frac{\|\mu - \nu\|_1} 2 = \sum_{x: \mu(x) > \nu(x)} (\mu(x) - \nu(x)) =  \sum_{x: \mu(x) > \nu(x)} \frac{(\mu(x) - \nu(x))^2}{\sqrt{\mu(x) + \nu(x)}} \sqrt{\mu(x) + \nu(x)} \underset{(*)}{\leq} \sqrt{2\Lambda(\mu, \nu)}.
\]
Here $(*)$ follows from Cauchy-Schwarz. To see the upper bound:
\[
\Lambda(\mu, \nu) = \sum_{x: \mu(x) > \nu(x)}\frac{(\mu(x) - \nu(x))^2}{\mu(x) + \nu(x)} = \sum_{x: \mu(x) > \nu(x)}\frac{(\mu(x) - \nu(x))^2}{\mu(x) - \nu(x)} \cdot \frac{\mu(x) - \nu(x)}{\mu(x) + \nu(x)} \underset{(**)}{\leq} \frac{\|\mu - \nu\|_1} 2.
\]
Here $(**)$ follows from $\frac{\mu(x) - \nu(x)}{\mu(x) + \nu(x)}\leq 1$.

\end{proof}

We will make use of the following:

\begin{proposition}[\cite{Yehudayoff20}]
\label{fact:trick1}
If $|\eta| \leq \sqrt{a(b+ \eta)}$ with $a,b\geq 0$, then $\eta \leq a + 2b$.
\end{proposition}

\begin{proof}
    If $a = 0$, then $\eta = 0$ and the upper-bound follows. Otherwise, if $\eta^2 - a \eta - a b \le 0$, then
    \[
    \eta \le \frac{a + \sqrt{a^2 + 4 a b}}{2} = \frac{a}{2} (1 + \sqrt{1 + 4 b/a}) \le \frac{a}{2}(1 + 1 + 4 b/a). \qedhere
    \]
\end{proof}

\lambdaloss*

\begin{proof}
We have:
\begin{align*}
    |\eta|=\left|\E_\mu(f(x)) - \E_\nu(f(x))\right| &= \left|\sum_x(\mu(x) - \nu(x))f(x)\right|\\
    &\leq \sum_x\left|(\mu(x) - \nu(x))\right|f(x)\tag{By triangle inequality}\\
    &= \sum_x\frac{|\mu(x) - \nu(x)|\sqrt{f(x)}}{\sqrt{\mu(x) + \nu(x)}} \sqrt{(\mu(x) + \nu(x)) f(x)}\\
    &\leq \sqrt{\sum_x\frac{(\mu(x) - \nu(x))^2f(x)}{\mu(x) + \nu(x)}} \sqrt{\sum_x(\mu(x) + \nu(x)) f(x)}\tag{By Cauchy-Schwarz}
\end{align*}

\noindent
Let us look at the first term without the square root:
\begin{align*}
    \sum_x\frac{(\mu(x) - \nu(x))^2f(x)}{\mu(x) + \nu(x)} &= \sum_{x: \mu(x) \geq \nu(x)}\frac{(\mu(x) - \nu(x))^2f(x)}{\mu(x) + \nu(x)} + \sum_{x: \mu(x) < \nu(x)}\frac{(\mu(x) - \nu(x))^2f(x)}{\mu(x) + \nu(x)}\\
    &\leq \Lambda(\mu,\nu)\cdot \max_x f(x) + \sum_{x: \mu(x) < \nu(x)}\frac{(\mu(x) - \nu(x))^2f(x)}{\mu(x) + \nu(x)}\\
    &\leq \Lambda(\mu,\nu)\cdot \max_x f(x) + \sum_{x}\frac{(\nu(x))^2}{\nu(x)}f(x) = \Lambda(\mu,\nu)\cdot \max_x f(x) + \E_\nu(f(x)).
\end{align*}
which implies:
\begin{align*}
    |\eta|=\left|\E_\mu(f(x)) - \E_\nu(f(x))\right| \leq \sqrt{\underbrace{(\Lambda(\mu,\nu)\cdot \max_x f(x) + \E_\nu(f(x)))}_{a} (\eta+ \underbrace{2\E_\nu(f(x)))}_b)}
\end{align*}
At this point, we can use \Cref{fact:trick1} to get the following:
\begin{align*}
    \E_\mu(f(x)) - \E_\nu(f(x)) \leq \Lambda(\mu,\nu)\cdot \max_x f(x) + 5\Pr_\nu(\cE),\\
    \text{i.e., }\E_\mu(f(x))\leq \Lambda(\mu,\nu)\cdot \max_x f(x) + 6\E_\nu(f(x)),
\end{align*}
which completes the proof.
\end{proof}

\convex*

\begin{proof}
To see this, it suffices to compute the Hessian of each term.
We have
\[
\Lambda(\mu,\nu) = \sum_{x} \begin{cases}
\frac{(\mu(x)-\nu(x))^2}{\mu(x)+\nu(x)} & \text{if } \mu(x) \ge \nu(x)\\
0 & \text{otherwise}.
\end{cases}
\]
So let us prove that each term is jointly convex in the values of $(\mu(x), \nu(x))$. We restrict our attention to the unit square. The function is zero, hence convex, inside the region where $\mu(x) \le \nu(x)$. Furthermore, it is (obviously) non-negative inside the region where $\mu(x) > \nu(x)$. It is also convex in this region, since the Hessian is positive semidefinite there:
 \[
 \nabla^2 \frac{(\mu(x) - \nu(x))^2}{\mu(x)+\nu(x)}  =
\begin{pmatrix}
 \frac{8 \nu(x)^2}{(\mu(x) +\nu(x) )^3} & -\frac{8 \mu(x)  \nu(x) }{(\mu(x) +\nu(x) )^3} \\
 -\frac{8 \mu(x)  \nu(x) }{(\mu(x) +\nu(x) )^3} & \frac{8 \mu(x)^2}{(\mu(x) +\nu(x) )^3}
\end{pmatrix}.
\]
(The matrix is PSD since the off-diagonal entry, in absolute value, does not exceed the geometric mean of the diagonal entries.) It follows that each term is convex, and hence $\Lambda(\mu,\nu)$ is convex, in $(\mu,\nu)$.
\end{proof}



\section{Proof of \texorpdfstring{\Cref{clm:nonmdirect-sum}}{Claim}} \label{sec:proof-direct-sum}

\begin{proof}
This proof is a verbatim reproduction of that in \cite{AssadiCK19}. For any $i \in [m]$, define $\bA^{<i} := \{A_{x_1},\ldots,A_{x_{i-1}}\}$, $\bB^{>i} := \{B_{x_{i+1}},\ldots,B_{x_{m}}\}$. Recall that 
the internal information cost of $\prot_\SI$ is $\IC{\prot_\SI}{\cD_\SI} := \mi{\rA}{\rProtSI \mid \rB} + \mi{\rB}{\rProtSI \mid \rA}$. In the following, we focus on bounding the first term. The second term can be bounded exactly
the same by symmetry.  

As $(\rI,\bA^{<\rI},\bB^{>\rI},\bC,\bD)$ is sampled via public randomness in $\prot_\SI$, by Proposition~\ref{prop:public-random},   
\begin{align*}
	\mi{\rA}{\rProtSI \mid \rB} &= \mi{\rA}{\rProtSI \mid \rB, \rI,\bA^{<\rI},\bB^{>\rI},\bC,\bD} \leq \mi{\rA}{\rProtSI \mid \rB, \rI,\bA^{<\rI},\bB^{>\rI}}. 
\end{align*}
The inequality is by Proposition~\ref{prop:info-decrease} as we now show $\rA \perp (\bC,\bD) \mid \rProtSI,\rB,\rI,\bA^{<\rI},\bB^{>\rI}$ (and hence conditioning on $(\bC,\bD)$ can only decrease the mutual information). 
This is because $\rA \perp (\bC,\bD) \mid \rB,\rI,\bA^{<\rI},\bB^{>\rI}$ by Observation~\ref{obs:distHPC} and $\rProtSI$ is transcript of a deterministic protocol plus $z_1,\ldots,z_j$ obtained deterministically and hence we can apply
Proposition~\ref{prop:hpc-rectangle}.  

Define a random variable $\rTheta \in \{0,1\}$ where $\rTheta = 1$ iff in Line~(\ref{line:AB}) of protocol $\prot_\SI$, we terminate the protocol. In other words $\rTheta = 1$ iff $x_i \in Z^{<j}$. Since $\rA \perp \rTheta \mid \rB, \rI,\bA^{<\rI},\bB^{>\rI}$, 
further conditioning on $\rTheta$ can only increase the mutual information term above by Proposition~\ref{prop:info-increase}, hence, 
\begin{align}
	\mi{\rA}{\rProtSI \mid \rB} &\leq \mi{\rA}{\rProtSI \mid \rB, \rI,\bA^{<\rI},\bB^{>\rI},\rTheta} \notag \\
	&= \frac{m-j}{m} \cdot \mi{\rA}{\rProtSI \mid \rB, \rI,\bA^{<\rI},\bB^{>\rI},\rTheta=0} + \frac{j}{m} \cdot \mi{\rA}{\rProtSI \mid \rB, \rI,\bA^{<\rI},\bB^{>\rI},\rTheta=1} \notag \\
	&\leq \frac{m-j}{m} \cdot \mi{\rA}{\rProtSI \mid \rB, \rI,\bA^{<\rI},\bB^{>\rI},\rTheta=0}, \label{eq:mi-ra}
\end{align}
since conditioned on $\rTheta=1$, the protocol $\rProtSI$ is simple some prefix of $\rZ^{<j}$ and is hence independent of the input $(\rA,\rB)$ and carries no information about $\rA$ (see \itfacts{info-zero}). 
We now further bound the RHS of Eq~(\ref{eq:mi-ra}). When $\rTheta = 0$, $\rProtSI = (\rZ^{<j} ,\rProt_1,\ldots,\rProt_j) = (\rE^{<j},\rProt_j)$. Hence, we can write, 
  \begin{align*}
	 \mi{\rA}{\rProtSI \mid \rB, \rI,\bA^{<\rI},\bB^{>\rI},\rTheta=0} &\leq \mi{\rA}{\rE_j,\rProt_j \mid \rB, \rI,\bA^{<\rI},\bB^{>\rI},\rTheta=0} \\
	&= \mi{\rA}{\rZ^{<j} \mid \rB, \rI,\bA^{<\rI},\bB^{>\rI},\rTheta=0} \\
	&\hspace{25pt} + \mi{\rA}{\rProt^{<j},\rProt_j \mid \rZ^{<j},\rB, \rI,\bA^{<\rI},\bB^{>\rI},\rTheta=0} \tag{by chain rule in \itfacts{chain-rule} and since $\rE_j = (\rProt^{<j},\rZ^{<j})$} \\
	&\leq \mi{\rA}{\rProt \mid \rZ^{<j},\rB, \rI,\bA^{<\rI},\bB^{>\rI},\rTheta=0},
\end{align*}
as $\rA \perp \rZ^{<j} \mid \rTheta=0$ (and other variables) and hence the first term is zero, and in the second term $\rProt$ contains $\rProt^{<j},\rProt_j$ (plus potentially other terms) and so having $\rProt$ in instead can only increase
the information. By further expanding the conditional information term above,
\begin{align*}
	 &\mi{\rA}{\rProtSI \mid \rB, \rI,\bA^{<\rI},\bB^{>\rI},\rTheta=0} \\
	 &\hspace{50pt}\leq \E_{(Z^{<j},i) \mid \rTheta=0} \bracket{\mi{\rA}{\rProt \mid \rB, \bA^{<i},\bB^{>i},\rI=i,\rZ^{<j}=Z^{<j},\rTheta=0}} \\
	&\hspace{50pt}= \E_{Z^{<j} \mid \rTheta=0} \bracket{\sum_{\substack{i=1 \\~i \notin Z^{<j}}}^{m} \frac{1}{m-j}  \mi{\rA_{x_i}}{\rProt \mid \rB_{x_i}, \bA^{<i},\bB^{>i},\rI=i,\rZ^{<j}=Z^{<j},\rTheta=0}} \tag{conditioned on $\rTheta=0$, $i$ is chosen
	uniformly at random from $Z^{<j}$; also $(\rA,\rB) = (\rA_{x_i},\rB_{x_i})$} \\
	&\hspace{50pt}= \E_{Z^{<j} \mid \rTheta=0} \bracket{\sum_{i \notin Z^{<j}} \frac{1}{m-j} \cdot \mi{\rA_{x_i}}{\rProt \mid \rB_{x_i}, \bA^{<i},\bB^{>i},\rZ^{<j}=Z^{<j},\rTheta=0}}
	\tag{we dropped the conditioning on $\rI=i$ as all remaining variables are independent of this event}  \\
	&\hspace{50pt}= \E_{Z^{<j} \mid \rTheta=0} \bracket{\sum_{i \notin Z^{<j}} \frac{1}{m-j} \cdot \mi{\rA_{x_i}}{\rProt \mid \bA^{<i},\bB,\rZ^{<j}=Z^{<j},\rTheta=0}}
	\tag{as $\rA_{x_i} \perp \bB^{<i} \mid \rB_{x_i},\bA^{<i}$ by Observation~\ref{obs:distHPC} and hence we can apply Proposition~\ref{prop:info-increase}} \\
	&\hspace{50pt}\leq \E_{Z^{<j} \mid \rTheta=0} \bracket{\sum_{i=1}^{m} \frac{1}{m-j} \cdot \mi{\rA_{x_i}}{\rProt \mid \bA^{<i},\bB,\rZ^{<j}=Z^{<j},\rTheta=0}} \\ 
	\tag{mutual information is non-negative by \itfacts{info-zero} and so we can add the terms in $Z^{<j}$  as well} \\
	&\hspace{50pt}= \E_{Z^{<j} \mid \rTheta=0} \bracket{\sum_{i=1}^{m} \frac{1}{m-j} \cdot \mi{\rA_{x_i}}{\rProt \mid \bA^{<i},\bB,\rZ^{<j}=Z^{<j},\rTheta=0}} \\
	&\hspace{50pt}= \frac{1}{m-j} \cdot \E_{Z^{<j} \mid \rTheta=0} \bracket{\mi{\bA}{\rProt \mid \bB,\rZ^{<j}=Z^{<j},\rTheta=0}} \tag{by chain rule in \itfacts{chain-rule}} \\
	&\hspace{50pt}= \frac{1}{m-j} \cdot \mi{\bA}{\rProt \mid \bB,\rZ^{<j},\rTheta=0} \\
	&\hspace{50pt}\leq \frac{1}{m-j} \cdot \paren{\mi{\bA}{\rProt \mid \bB,\rTheta=0} + \en{\rZ^{<j}}} \\
	&\hspace{50pt}= \frac{1}{m-j} \cdot \paren{\mi{\bA}{\rProt \mid \bB} + \en{\rZ^{<j}}} \tag{transcript of the protocol $\prot_\bhpc$ (namely $\rProt$) on input $(\bA,\bB)$ is independent of $\rTheta$} \\
	&\hspace{50pt}\leq \frac{1}{m-j} \cdot \paren{\en{\rProt} + \en{\rZ^{<j}}} \leq \frac{\CC{\prot_\bhpc}{} }{m-j}   + \frac{j \cdot \log{m}}{m-j}. \tag{by sub-additivity of entropy (\itfacts{sub-additivity}) and \itfacts{uniform}}
\end{align*}

By plugging in this bound in Eq~(\ref{eq:mi-ra}), we have that,
\begin{align*}
	\mi{\rA}{\rProtSI \mid \rB} &\leq \frac{m-j}{m} \cdot \paren{\frac{\CC{\prot_\bhpc}{} }{m-j}   + \frac{j \cdot \log{m}}{m-j}} = \frac{\CC{\prot_\bhpc}{} }{m}   + \frac{j \cdot \log{m}}{m}. 
\end{align*}
By symmetry, we can also prove the same bound on $\mi{\rB}{\rProtSI \mid \rA}$. As such, we have,
\begin{align*}
	\mi{\rA}{\rProtSI \mid \rB} + \mi{\rB}{\rProtSI \mid \rA} \leq 2 \cdot \paren{\frac{\CC{\prot_\bhpc}{} }{m}   + \frac{j \cdot \log{m}}{m}}.
\end{align*}

We shall note that strictly speaking the factor $2$ above is not needed (similar to the proof of Proposition~\ref{prop:cc-ic}) but as this factor is anyway suppressed through O-notation later in the proof, the above bound
suffices for our purpose. %\qed\; %\Cref{clm:direct-sum}

\end{proof}


\section{Proof of \texorpdfstring{\Cref{lem:redn-aux-constr}}{Claim}}\label{app:proof-auxnode-redn}
    Consider the graph $G'=(V',E')$. Set $d=6mr+3m$. We add $d$ auxiliary nodes in $\Vaux$. Since the target degree of each node outside $S$ is at most $d$, and the nodes in $S$ need at most $6r+3m\leq d$ edges to reach target degree, the $d$ auxiliary nodes suffice to increase their degree. Thus, the total number of vertices in $G$ is $|V'|+d=\Theta(mr)$. 
    
    We now describe the construction. For each node $v\in V'$, define its deficiency $\defc(v)$ as the difference between its target degree mentioned above and its current degree in $G'$. Define the total deficiency $D:= \sum_{v\in V'} \defc(v)$.  Fix arbitrary orderings $v_1,v_2,\ldots v_{|V'|}$ of the nodes in $V'$ and $u_1,\ldots,u_d$ of the nodes in $\Vaux$. From each $v_i$, we join edges to $\defc(v_i)$ many $u_j$'s following the ordering. The first neighbor of $v_{i+1}$ is taken to be the node that is next to the last neighbor of $v_i$ in the ordering of the $u_j$'s. We wrap around to $u_1$ after $u_d$. Hence, the degree of two vertices in $\Vaux$ can differ by at most $1$. We check the min-degree of a node in $\Vaux$. If it's $\geq d+6r+3$, then we are done. Otherwise, suppose it's $d+6r+3-x$ for some $x>0$. Then we add $x$ disjoint matchings of size $d/2$ in $\Vaux$ to increase the degree of each node by $x$, i.e., to ensure that each node in $\Vaux$ has degree $\geq d+6r+3$.  

    It remains to prove the ``furthermore'' part. For this, we check the maximum and minimum deficiency of each node.

    Observe that in $G'$, each node in layer $0$ has degree at least $5$ (it's adjacent to the other $2$ nodes in its triple and to the $3$ nodes in $S$) and at most $2m+5$ (for each $i\in [m]$, it is adjacent to at most $2$ nodes that represent $i$ in layer $1$). Hence, the deficiency of any such node is in $[d-2m-5, d-5]$. 
    
    For any node in a layer $\ell\in [2r-1]$, it is adjacent to at least $8$ nodes (the other $2$ nodes in its triple, the $3$ nodes in $S$, and the $3$ nodes representing the same element in its ``paired layer'', i.e., layer $\ell+1$ if $\ell$ is odd, or layer $\ell-1$ if it's even) and at most $2m+8$ nodes (for each $i\in [m]$, it is adjacent to at most $2$ nodes that represent $i$ in layer $\ell-1$ if $\ell$ is odd, or layer $\ell+1$ if it's even). Thus, the deficiency of any such node is in $[d-2m-8, d-8]$.

    Any node in layer $2r$ representing $x_i$ with $b(x_i)=0$ has degree exactly $5$ (the $2$ nodes in its triple and the $3$ nodes representing the same element in layer $2r-1$). Finally, any other nodes in layer $2r$ has degree exactly $8$ (same as other nodes in layer $2r$ with the addition of the $3$ nodes in $S$). Hence, the deficiency of any such node is in $[d-8, d-5]$. 

    Therefore, the total deficiency $D$ is at least $(d-2m-8)\cdot |V_L|= (d-2m-8)\cdot d \geq 9d^2/10$ for large enough $m$. These edges are received almost equally by the $d$ nodes in $|\Vaux|$. Thus, the min-degree in $|\Vaux|$ after the addition of these edges is at least $9d/10$. Thus the value of $x$ is at most $d+6r+3-9d/10\leq d-3$ (again for large enough $m$). The $x$ matchings can induce degree at most $x\leq d-3$ in the subgraph $G[\Vaux]$, and hence we are done with the furthermore part.    



%\clearpage
\section{Background on Information Theory}\label{prelim-sec:info}

We now briefly introduce some definitions and facts from information theory that are needed in this thesis. We refer the interested reader to the text by Cover and Thomas~\cite{CoverT06} for an excellent introduction to this field, 
and the proofs of the statements used in this Appendix. 

For a random variable $\rA$, we use $\supp{\rA}$ to denote the support of $\rA$ and $\distribution{\rA}$ to denote its distribution. 
When it is clear from the context, we may abuse the notation and use $\rA$ directly instead of $\distribution{\rA}$, for example, write 
$A \sim \rA$ to mean $A \sim \distribution{\rA}$, i.e., $A$ is sampled from the distribution of random variable $\rA$. 


We denote the \emph{Shannon Entropy} of a random variable $\rA$ by
$\en{\rA}$, which is defined as: 
\begin{align}
	\en{\rA} := \sum_{A \in \supp{\rA}} \Pr\paren{\rA = A} \cdot \log{\paren{1/\Pr\paren{\rA = A}}} \label{eq:entropy}
\end{align} 

\noindent
The \emph{conditional entropy} of $\rA$ conditioned on $\rB$ is denoted by $\en{\rA \mid \rB}$ and defined as:
\begin{align}
\en{\rA \mid \rB} := \E_{B \sim \rB} \bracket{\en{\rA \mid \rB = B}}, \label{eq:cond-entropy}
\end{align}
where 
$\en{\rA \mid \rB = B}$ is defined in a standard way by using the distribution of $\rA$ conditioned on the event $\rB = B$ in Eq~(\ref{eq:entropy}). 

\noindent
Recall that the binary entropy $H(p)$ is the entropy of a two-valued random
variable that takes one value with probability $p$.
We need the following upper bound on the binary entropy.

\begin{fact}\label{fact:ent-ub}
  The binary entropy $H_2(p) \leq 2 \sqrt{ p \cdot (1 - p) }$.
\end{fact}


 
The \emph{mutual information} of two random variables $\rA$ and $\rB$ is denoted by
$\mi{\rA}{\rB}$ and is defined as:
\begin{align}
\mi{\rA}{\rB} := \en{A} - \en{A \mid  B} = \en{B} - \en{B \mid  A}. \label{eq:mi}
\end{align}
\noindent
The \emph{conditional mutual information} $\mi{\rA}{\rB \mid \rC}$ is $\en{\rA \mid \rC} - \en{\rA \mid \rB,\rC}$ and hence by linearity of expectation:
\begin{align}
	\mi{\rA}{\rB \mid \rC} = \E_{C \sim \rC} \bracket{\mi{\rA}{\rB \mid \rC = C}}. \label{eq:cond-mi}
\end{align} 



\subsection{Useful Properties of Entropy and Mutual Information}\label{sec:prop-en-mi}

We shall use the following basic properties of entropy and mutual information throughout. 


\begin{fact}\label{fact:it-facts}
  Let $\rA$, $\rB$, $\rC$, and $\rD$ be four (possibly correlated) random variables.
   \begin{enumerate}
  \item \label{part:uniform} $0 \leq \en{\rA} \leq \log|\supp{\rA}|$. The right equality holds
    iff $\distribution{\rA}$ is uniform.
  \item \label{part:info-zero} $\mi{\rA}{\rB \mid \rC} \geq 0$. The equality holds iff $\rA$ and
    $\rB$ are \emph{independent} conditioned on $\rC$.
    \item \label{part:info-atmost-rv} $\mi{\rA}{\rB \mid \rC} \leq \en{\rA}$.
  \item \label{part:cond-reduce} \emph{Conditioning on a random variable reduces entropy}:
    $\en{\rA \mid \rB,\rC} \leq \en{\rA \mid  \rB}$.  The equality holds iff $\rA \perp \rC \mid \rB$.
    \item \label{part:sub-additivity} \emph{Subadditivity of entropy}: $\en{\rA,\rB \mid \rC}
    \leq \en{\rA \mid C} + \en{\rB \mid  \rC}$.
   \item \label{part:ent-chain-rule} \emph{Chain rule for entropy}: $\en{\rA,\rB \mid \rC} = \en{\rA \mid \rC} + \en{\rB \mid \rC,\rA}$.
  \item \label{part:chain-rule} \emph{Chain rule for mutual information}: $\mi{\rA,\rB}{\rC \mid \rD} = \mi{\rA}{\rC \mid \rD} + \mi{\rB}{\rC \mid  \rA,\rD}$.
  \item \label{part:data-processing} \emph{Data processing inequality}: for a function $f(\rA)$ of $\rA$, $\mi{f(\rA)}{\rB \mid \rC} \leq \mi{\rA}{\rB \mid \rC}$. 
   \end{enumerate}
\end{fact}


\noindent
We also use the following two standard propositions, regarding the effect of conditioning on mutual information.

\begin{proposition}\label{prop:info-increase}
  For random variables $\rA, \rB, \rC, \rD$, if $\rA \perp \rD \mid \rC$, then, 
  \[\mi{\rA}{\rB \mid \rC} \leq \mi{\rA}{\rB \mid  \rC,  \rD}.\]
\end{proposition}
 \begin{proof}
  Since $\rA$ and $\rD$ are independent conditioned on $\rC$, by
  \itfacts{cond-reduce}, $\HH(\rA \mid  \rC) = \HH(\rA \mid \rC, \rD)$ and $\HH(\rA \mid  \rC, \rB) \ge \HH(\rA \mid  \rC, \rB, \rD)$.  We have,
	 \begin{align*}
	  \mi{\rA}{\rB \mid  \rC} &= \HH(\rA \mid \rC) - \HH(\rA \mid \rC, \rB) = \HH(\rA \mid  \rC, \rD) - \HH(\rA \mid \rC, \rB) \\
	  &\leq \HH(\rA \mid \rC, \rD) - \HH(\rA \mid \rC, \rB, \rD) = \mi{\rA}{\rB \mid \rC, \rD}. \qed
	\end{align*}
	
\end{proof}

\begin{proposition}\label{prop:info-decrease}
  For random variables $\rA, \rB, \rC,\rD$, if $ \rA \perp \rD \mid \rB,\rC$, then, 
  \[\mi{\rA}{\rB \mid \rC} \geq \mi{\rA}{\rB \mid \rC, \rD}.\]
\end{proposition}
 \begin{proof}
 Since $\rA \perp \rD \mid \rB,\rC$, by \itfacts{cond-reduce}, $\HH(\rA \mid \rB,\rC) = \HH(\rA \mid \rB,\rC,\rD)$. Moreover, since conditioning can only reduce the entropy (again by \itfacts{cond-reduce}), 
  \begin{align*}
 	\mi{\rA}{\rB \mid  \rC} &= \HH(\rA \mid \rC) - \HH(\rA \mid \rB,\rC) \geq \HH(\rA \mid \rD,\rC) - \HH(\rA \mid \rB,\rC) \\
	&= \HH(\rA \mid \rD,\rC) - \HH(\rA \mid \rB,\rC,\rD) = \mi{\rA}{\rB \mid \rC,\rD}. \qed
 \end{align*}

\end{proof}


\subsection{Background on Communication and Information Complexity}\label{sec:cc-ic}

\paragraph{Communication complexity.} We briefly review the standard definitions of the two-party communication model of Yao~\cite{Yao79}.  See the text by Kushilevitz and Nisan~\cite{KushilevitzN97} for 
an extensive overview of communication complexity. In Section~\ref{sec:lower-bound-hpc}, we also use a standard generalization of this model to allow for more than two players, but we defer
the necessary definitions to that section. 

Let $P: \mathcal{X} \times \mathcal{Y} \rightarrow \mathcal{Z}$ be a relation.  Alice receives an input $X
\in \mathcal{X}$ and Bob receives $Y \in \mathcal{Y}$, where $(X,Y)$ are chosen from a
joint distribution $\cD$ over $\mathcal{X} \times \mathcal{Y}$. We allow players to have access to both public and private randomness. 
They communicate with each other by exchanging messages such that each message
depends only on the private input and random bits of the player sending the message, and the already communicated messages plus the public randomness. 
At the end, one of the players need to output an answer  $Z$ such that $Z \in P(X,Y)$.  

We use $\prot$ to denote a protocol used by the players. We always assume that the protocol $\prot$ can be randomized (using both public and
private randomness), \emph{even against a prior distribution $\cD$ of inputs}. For any
$0 < \delta < 1$, we say $\prot$ is a $\delta$-error protocol for $P$ over a distribution
$\cD$, if the probability that for an input $(X,Y)$, $\prot$ outputs some $Z$ where $Z \notin P(X,Y)$ is at most
$\delta$ (the probability is taken over the randomness of \emph{both} the distribution and the protocol).

\begin{definition}[Communication cost]
  The \emph{communication cost} of a protocol $\prot$ on an input
  distribution $\cD$, denoted by $\CC{\prot}{\cD}$, is the \emph{worst-case} bit-length of the transcript
  communicated between Alice and Bob in the protocol $\prot$, when the inputs are chosen from $\cD$.

\end{definition}

Communication complexity of a problem $P$ is defined as the minimum communication cost of a protocol $\prot$ that solves $P$ on every distribution $\cD$ with probability at least $2/3$. 



\paragraph{Information complexity.} There are several possible definitions of information
cost of a communication prtocol that have been considered depending on the application (see, e.g.,~\cite{ChakrabartiSWY01,Bar-YossefJKS02,BarakBCR10,BravermanR11,BravermanEOPV13}).  
We use the notion of \emph{internal information cost}~\cite{BarakBCR10} that measures the average amount of information each player learns about the input of the other player
by observing the transcript of the protocol. %Formally, 
\begin{definition}[Information cost]
  Consider an input distribution $\cD$ and a protocol $\prot$. Let $(\rX,\rY) \sim \cD$ denote the random variables for the input of Alice and Bob and $\rProt$ be the 
  the random variable for the transcript of the protocol \emph{concatenated} with the public randomness $\rR$ used by $\prot$. 
  The \emph{(internal) information cost} of $\prot$ with respect to
  $\cD$ is $\ICost{\prot}{\cD}:=\mii{\rProt}{\rX \mid \rY}{\cD} + \mii{\rProt}{\rY \mid \rX}{\cD}$. 

\end{definition}

One can also define information complexity of a problem $P$ similar to communication complexity with respect to the information cost. However, we avoid presenting this definition formally due to some subtle technical issues
 that need to be addressed which lead to multiple different but similar-in-spirit definitions. As such, we state our results directly in terms of information cost. 

%%
Note that any public coin protocol is a distribution over private coins protocols, run by
first using public randomness to sample a random string $\rR=R$ and then running the
corresponding private coin protocol $\prot^R$. We also use $\rProt^R$ to denote the transcript of the protocol $\prot^R$. 
We have the following standard proposition. 
\begin{proposition}\label{prop:public-random}
	 For any distribution $\cD$ and any protocol $\prot$ with public randomness $\bR$, $$\ICost{\prot}{\cD} = \mii{\rProt}{\rX \mid \rY, \rR}{\cD} + \mii{\rProt}{\rY \mid \rX, \rR}{\cD} = \E_{R \sim \rR}\bracket{\ICost{\prot^R}{\cD}}.$$ 
\end{proposition}
\begin{proof} By definition of internal information cost, 
\begin{align*}
	\ICost{\prot}{\cD} &= \mii{\rProt}{\rX \mid \rY}{\cD} + \mii{\rProt }{\rY \mid \rX}{\cD} = \mi{\rProt,\rR}{\rX \mid \rY} + \mi{\rProt,\rR}{\rY \mid \rX} \tag{$\Prot$ denotes the transcript and the public randomness} \\
	&= \mi{\rR}{\rX \mid \rY}  + \mi{\rProt}{\rX \mid \rY , \rR} + \mi{\rR}{\rY \mid \rX}  + \mi{\rProt}{\rY \mid \rX , \rR}  \tag{chain rule of mutual information, \itfacts{chain-rule}} \\
	&= \mi{\rProt}{\rX \mid \rY , \rR} + \mi{\rProt}{\rY \mid \rX , \rR} \tag{$\mi{\rR}{\rX \mid \rY}  = \mi{\rR}{\rY \mid \rX}  = 0$ since $\rR \perp \rX, \rY$ and \itfacts{info-zero}} \\
	&= \E_{R \sim \rR}\bracket{\mi{\rProt}{\rX \mid \rY , \rR=R} + \mi{\rProt}{\rY \mid \rX , \rR=R}} = \E_{R \sim \rR}\bracket{\ICost{\prot^R}{\cD}}, 
\end{align*}
concluding the proof.
\end{proof}


The following well-known proposition relates communication cost and information cost. 
\begin{proposition}[cf.~\cite{BravermanR11}]\label{prop:cc-ic}
  For any distribution $\cD$ and any protocol $\prot$: $\ICost{\prot}{\cD} \leq \CC{\prot}{\cD}$.
\end{proposition}
\begin{proof}
Let us assume first that $\prot$ only uses private randomness and thus $\rProt$ only contain the transcript.  
For any $b \in [\CC{\prot}{\cD}]$, we define $\Prot_b$ to be the $b$-th bit of the transcript. We have, 
\begin{align*}
	 \ICost{\prot}{\cD} &= \mi{\rProt}{\rX \mid \rY}+ \mi{\rProt }{\rY \mid \rX} \\
	 &= \sum_{b=1}^{\CC{\prot}{\cD}} \mi{\rProt_b}{\rX \mid \rProt^{<b},\rY} + \mi{\rProt_b}{\rY \mid \rProt^{<b},\rX} \tag{by chain rule of mutual information in \itfacts{chain-rule}}\\
	 &= \sum_{b=1}^{\CC{\prot}{\cD}} \E_{\Prot^{<b}}\bracket{\mi{\rProt_b}{\rX \mid \rProt^{<b} = \Prot^{<b},\rY} + \mi{\rProt_b}{\rY \mid \rProt^{<b}= \Prot^{<b},\rX}}.
\end{align*}
Consider each term in the RHS above. By conditioning on $\Prot^{<b}$, the player that transmit $\rProt_b$ would become fix. If this player is Alice, then $\mi{\rProt_b}{\rY \mid \rProt^{<b}= \Prot^{<b},\rX} = 0$, because
$\rProt_b$ is only a function of $(\rProt^{<b},\rX)$ in this case; similarly, if this player is Bob, then $\mi{\rProt_b}{\rX \mid \rProt^{<b}= \Prot^{<b},\rY} = 0$. Moreover, 
$\mi{\rProt_b}{\rX \mid \rProt^{<b} = \Prot^{<b},\rY} \leq \en{\rProt_b}  \leq 1$ and similarly $\mi{\rProt_b}{\rY \mid \rProt^{<b}= \Prot^{<b},\rX} \leq 1$. As such, the above term can be upper bounded 
by $\CC{\prot}{\cD}$. To finalize the proof, note that by Proposition~\ref{prop:public-random}, for any public-coin protocol $\prot$, $\ICost{\prot}{\cD} = \E_{R \sim \rR}\bracket{\ICost{\prot^R}{\cD}} \leq \E_{R \sim \rR}\bracket{\CC{\prot^R}{\cD}} \leq \CC{\prot}{\cD}$, where
the first inequality is by the first part of the argument.  
\end{proof}

Proposition~\ref{prop:cc-ic} provides a convinent way of proving communication complexity lower bounds by lower bounding information cost of any protocol. 


\subsection*{Rectangle Property of Communication Protocols}\label{sec:statstics}
We conclude this section by mentioning some basic properties of communication protocols. For any protocol $\prot$ and inputs $x \in \mathcal{X}$ and $y \in \mathcal{Y}$, we define $\Prot_{x,y}$ as the transcript
of the protocol conditioned on the input $x$ to Alice and input $y$ to Bob. Note that for randomized protocols, $\Prot_{x,y}$ is a random variable which we denote by $\rProt_{x,y}$. 

The following is referred to as the rectangle property of deterministic protocols. 
\begin{fact}[Rectangle property]\label{fact:rectangle}
	For any deterministic protocol $\prot$ and inputs $x,x' \in \mathcal{X}$ to Alice and $y,y' \in \mathcal{Y}$ to Bob, if $\Prot_{x,y} = \Prot_{x',y'}$, then $\Prot_{x,y'} = \Prot_{x',y}$. 
\end{fact}
\noindent
Fact~\ref{fact:rectangle} implies that the set of inputs consistent with any transcript $\Prot_{x,y}$ of a deterministic protocol forms a combinatorial rectangle. 
One can also extend the rectangle property of deterministic protocols to randomized protocols using the following fact. 
\begin{fact}[Cut-and-paste property; cf.~\cite{Bar-YossefJKS02}]\label{fact:r-rectangle}
	For any randomized protocol $\prot$ and inputs $x,x' \in \mathcal{X}$ to Alice and $y,y' \in \mathcal{Y}$ to Bob, $\hd{\rProt_{x,y}}{\rProt_{x',y'}} = \hd{\rProt_{x,y'}}{\rProt_{x',y}}$. 
\end{fact}
